\section{Background and related work}\label{sec:background}

\subsection{Design-stage evaluation of lifecycle properties}

Systematic design methodology treats the evaluation of alternatives against lifecycle criteria as a core activity of the conceptual and embodiment phases \citep{Pahl2007, Tomiyama2009}, and the Design-for-X programme has turned a growing list of such criteria into operational methods \citep{Holt2010}. Reliability is the criterion closest to the present work. Failure mode and effects analysis remains the dominant design-stage instrument; recent work in this journal has re-founded it on design theory \citep{Cabanes2021}, extended it with causal reasoning over failure data \citep{Liu2026afgnn}, embedded it in frameworks that demonstrate high reliability in regulated products \citep{Booth2026}, and coupled it to probabilistic models of early design risk \citep{Truong2026} and to machine-learning surrogates of degradation in robust design \citep{Hassani2026}. What FMEA-type methods record about \emph{detection} is, however, an expert rating: the ``D'' column of a DFMEA is assigned, not measured, and it is assigned for one architecture at a time. The method proposed here can be read as a way of measuring that column across alternatives.

The comparison of architectural alternatives is itself a design-methodological topic. \citet{Eaton2026} show that the technical measures omitted from an alternative-selection problem change which alternative is chosen, and \citet{Park2019} analyse how the topology of a product family constrains its evolution. Diagnosability is, we argue, a routinely omitted technical measure; and because monitoring designs are increasingly expected to serve a family of variants rather than one product \citep{Simpson2004}, whether a monitoring design transfers across alternatives is a family-level question that the selection method should answer.

\subsection{Diagnosability, testability and sensor placement}

Fault detection and isolation (FDI) has a mature theory of diagnosability for systems described by models. Structural analysis determines which faults can be detected and isolated from which sensor sets without numerical parameter values \citep{Blanke2016, Krysander2008, Frisk2009}, and has been applied to electric-vehicle drives with a diagnosability index that ranks candidate residual generators \citep{Zhang2015}. In electronics, testability analysis and multi-signal flow graphs answer the analogous question for test points \citep{Simpson1994, Deb1994}. \citet{Basseville2001} gives statistical definitions of detectability and isolability for parametric models, and the prognostics community has developed sensor-selection criteria and performance metrics for health management \citep{Cheng2010, Saxena2008}. Model-based approaches \citep{Isermann2005, Venkatasubramanian2003} require a model whose structure is known and trusted; data-driven fault diagnosis \citep{Lei2020} requires data from the very system being designed, and its transfer across systems is an active problem \citep{Pan2010}.

The method proposed here occupies the space between these. It does not require a model structure, only trials with a known fault extent, which at the design stage can be generated by simulation and at the test stage by prototypes or benchmarks; it does not require a trained detector, only observable features; and it produces a quantity---the SDR---that is comparable across features, operating points and alternatives because it is normalised by the healthy system's own non-stationarity rather than by a physical unit. It is closer in spirit to change-detection statistics \citep{Basseville1993} than to residual generation, and it deliberately stops at detectability: isolation and prognosis are downstream questions that presuppose it.

\subsection{Condition monitoring of electrical machines}

Inter-turn short-circuits in stator windings are among the most studied incipient faults in electrical machines \citep{Gandhi2011, RieraGuasp2015}. Motor-current signature analysis \citep{Thomson2001}, negative-sequence current and impedance methods \citep{Kliman1996, Lee2003}, transient models of turn faults \citep{Tallam2002} and the review literature \citep{Nandi2005, Bellini2008} have established the physical signatures: an asymmetry of the winding produces negative-sequence components in the terminal quantities and a component at twice the fundamental frequency in the synchronous reference frame. Which of these appears where depends on whether the machine is fed from the grid or from a current-controlled converter, because a current controller acts against the asymmetry it can see. This observation---well known to drive specialists---is the physical basis of one of our design findings, but it has not been turned into a design-stage measure. Machine-learning detectors, including our own recent work on multi-fault detection in PMSM drives, on permanent-magnet faults in offshore generators identified from finite-element simulations by transfer learning, and on onset-aware diagnosis \citep{Canseven2026iet, Canseven2026machines, Sadik2026}, are typically developed and evaluated for one machine; the transferability experiments in this paper are the design-level counterpart of that line of work. Digital twins are the natural vehicle for the design-stage instantiation of the method \citep{Tao2019}; how such twins should themselves be designed is an open question in this journal \citep{Tipuric2026, Tadeja2026}.

\subsection{Generator topology as a wind-turbine design decision}

The generator concept of a large wind turbine---direct-drive or geared, permanent-magnet, induction or wound-field, full or partial power conversion---is selected early on cost, mass, efficiency and grid-compliance grounds \citep{Polinder2006, Polinder2013}. Reliability enters the trade mostly through failure-rate statistics of subassemblies \citep{Tavner2007, Carroll2016}, not through the observability of the failures. Modern machine design optimisation has become highly effective at the electromagnetic and thermal objectives \citep{Bramerdorfer2018, Canseven2024tia, Canseven2025tec}; the present method adds an objective that has so far been missing from that toolbox.
